%==============================================================================
%  STAGING FRAGMENT --- skew-brace decomposition obstruction [Omega]
%  ASSEMBLE FROM:
%    proofs/2026-09-25-bilinear-skew-brace-h2-identification.tex  [peer-reviewed]
%    proofs/2026-09-24-partial-pruning-bilinear-obstruction.tex   [proved]
%    proofs/2026-09-24-heap-transport-refuted-localization.tex    [proved]
%    proofs/2026-09-25-nu-variation-direct-factor.tex             [proved]
%    proofs/2026-09-25-nontrivial-kernel-orthogonality.tex        [proved]
%    proofs/2026-09-26-nu-omega-product-theorem.md                [proved]
%    proofs/2026-09-26-nu-omega-independence-corner-and-mechanism.md [computed]
%  STATUS 2026-09-26: standalone chapter; NOT yet \input into dossier.tex
%  (pending Neil's steer on whether this becomes its own pillar chapter).
%
%  Self-contained macro guard so the fragment compiles both alone (via the
%  _standalone wrapper) and after later \input into the dossier preamble.
%==============================================================================
\providecommand{\Aut}{\mathrm{Aut}}
\providecommand{\Soc}{\mathrm{Soc}}
\providecommand{\Sb}{\mathrm{Sb}}
\providecommand{\Gp}{\mathrm{Gp}}
\providecommand{\Ext}{\mathrm{Ext}}
\providecommand{\Hom}{\mathrm{Hom}}
\providecommand{\id}{\mathrm{id}}
\providecommand{\Z}{\mathbb{Z}}
\providecommand{\D}{\mathcal{D}}
\providecommand{\Hcoh}{H^2}
\newcommand{\bic}{\bowtie}          % bicrossed (Zappa--Szep) product
\newcommand{\ovtau}{\bar\tau}
\newcommand{\Om}{[\Omega]}
\providecommand{\ract}{\triangleright}
% NB on merge into dossier.tex: this fragment cites \cite{ferri-qsb} (already in
% the dossier bib) and \cite{RatheeYadavExt} (arXiv:2601.12371, deep-read), which
% must be ADDED to the dossier bibliography. Suggested bibitem:
%   \bibitem{RatheeYadavExt} N.~Rathee, M.~K.~Yadav, \emph{Skew brace extensions,
%     second cohomology and complements}, arXiv:2601.12371, 2026.
% CITATION TODO (browse session): the earlier Rathee-Yadav cohomology paper
% "Cohomology, extensions and automorphisms of skew braces" (origin of the
% beta=d theta convention) is only agent-summary in sources.json, BELOW the
% deep-read citation floor -- so it is deliberately NOT cited here. If a precise
% citation to it is wanted, deep-read it first and upgrade its sources entry.

\chapter{The decomposition obstruction of a skew brace: a bilinear class, independent of the internal residue}
\label{ch:skewbrace-omega}

A skew brace is the algebraic shadow of a set-theoretic solution of the
Yang--Baxter equation, and the recurring engineering question about it is the
same one that runs through this whole dossier: \emph{when does a composite come
apart into its pieces?} Concretely, given a skew brace $G$ and an ideal
$D\trianglelefteq G$, when is $G$ an internal bicrossed (Zappa--Sz\'ep) product
$(G/D)\bic D$? This chapter answers three questions about the obstruction to that
decomposition. It \emph{names} the obstruction (a class in Rathee--Yadav skew-brace
cohomology, \textsc{[peer-reviewed]}); it \emph{measures} the obstruction (a
genuinely \emph{bilinear} class, strictly finer than the two underlying group
Schreier classes, \textsc{[proved]}); and it \emph{isolates} the obstruction from
the one datum most easily confused with it --- the internal relabelling residue
$[\nu]$ --- by proving they are independent coordinates (\textsc{[proved]}, with
one surjectivity residual \textsc{[computed]}).

The result answers, in group-theoretic language, Ferri's Remark~3.20 on when a
quiver skew brace splits off an ideal \cite{ferri-qsb}, and it supplies the
skew-brace / Yang--Baxter instance of the dossier's decomposition programme. A
word of warning before we begin, because the two are so easily conflated: the
\emph{existence} axis --- ``does $G$ admit a Zappa--Sz\'ep factorisation at all''
--- is the subject of Chapter~\ref{ch:zs} and is governed by the group-cohomological
class $[\omega]$. The \emph{decomposition} axis of this chapter --- ``does $G$
factor \emph{through a prescribed ideal} $D$'' --- is strictly finer, and its
obstruction $\Om$ lives in a two-operation cohomology that no single group class
can see. Keeping the two apart is the point.

%------------------------------------------------------------------------------
\section{The three theorems}\label{sec:omega-intro}

Throughout, $(G,+,\circ)$ is a skew left brace: $(G,+)$ and $(G,\circ)$ are
groups sharing the identity $0$, related by
$a\circ(b+c)=(a\circ b)-a+(a\circ c)$, and we write
$\lambda_a(b):=-a+(a\circ b)$, so $a\circ b=a+\lambda_a(b)$ and
$\lambda\colon(G,\circ)\to\Aut(G,+)$ is a homomorphism. An \emph{ideal}
$D\trianglelefteq G$ is a normal subgroup for both operations with
$\lambda_a(D)=D$; it yields a short exact sequence of skew braces
\[
   0 \longrightarrow D \longrightarrow G
     \xrightarrow{\ \pi\ } H:=G/D \longrightarrow 0 .
\]
We assume $(D,+)$ and $(D,\circ)$ are both abelian --- the setting in which the
extension is classified by a second cohomology --- and, where a citation to
Rathee--Yadav's theory is invoked directly, the standing hypothesis (SH) that $D$
is a \emph{trivial} brace ($a\circ b=a+b$ on $D$). A \emph{decomposition of $G$
through $D$} is an isomorphism $G\cong (G/D)\bic D$ onto an internal bicrossed
product; equivalently (see \S\ref{sec:omega-identify}) a sub-skew-brace complement
to $D$.

Fix a normalised section $s\colon H\to G$ (so $s(0)=0$). It determines the
\emph{good triplet}
\[
  \nu_h := \lambda_{s(h)}\!\mid_D,\qquad
  \mu_h(x) := -s(h)+x+s(h),\qquad
  \sigma_h(x) := s(h)^{-1}\circ x\circ s(h),
\]
of maps $D\to D$, together with the additive and circle factor sets
\[
  \beta(h_1,h_2) := -s(h_1{+}h_2)+s(h_1)+s(h_2)\in D,\qquad
  \ovtau(h_1,h_2) := s(h_1{\circ}h_2)^{-1}\circ s(h_1)\circ s(h_2)\in D.
\]
The three theorems of the chapter, stated informally now and precisely in the body:

\begin{theorem}[Identification, \textsc{[peer-reviewed]}; \S\ref{sec:omega-identify}]
\label{thm:A-informal}
The obstruction to decomposing $G$ through $D$ is a single class
$\Om\in\Hcoh_{\Sb}(H,D)$ in Rathee--Yadav skew-brace cohomology, and
$\Om=0$ if and only if $G\cong(G/D)\bic D$, i.e.\ $D$ has a sub-skew-brace
complement. No coprimality hypothesis is needed.
\end{theorem}

\begin{theorem}[Bilinearity, \textsc{[proved]}; \S\ref{sec:omega-bilinear}]
\label{thm:B-informal}
$\Om$ is a genuinely \emph{bilinear} class: the two forgetful maps to the
additive and circle group Schreier classes,
$\varphi_+(\Om)=[\beta]$ and $\varphi_\circ(\Om)=[\ovtau]$, can both vanish while
$\Om\neq0$. Vanishing of \emph{both} single-operation obstructions does
\emph{not} imply decomposition. The obstruction is base-localised: it is decided
on a single star of the ambient groupoid, and closure there propagates to all
stars.
\end{theorem}

\begin{theorem}[Independence, \textsc{[proved]}; \S\S\ref{sec:omega-residue}--\ref{sec:omega-product}]
\label{thm:C-informal}
The internal relabelling residue $[\nu]\in\prod_h\Aut(D,+)/L$ (where
$L:=\operatorname{im}(\lambda\!\mid_D)$) is an \emph{orthogonal direct factor} of
the section-independent data, independent of $\Om$. Fixing $(H\text{-brace},\mu,
\sigma,L)$, the equivalence classes of extensions with a given realised residue
form a torsor under a residue-\emph{independent} obstruction group $\Omega^0$;
consequently
\(
   \bigl([\nu],\Om\bigr)\colon \Ext \xrightarrow{\ \sim\ } R\times\Omega^0
\)
is a bijection --- \emph{every realised residue co-occurs with every value of the
obstruction.} That every coupling-admissible residue is realised (surjectivity of
the residue map) is \textsc{[computed]}, not \textsc{[proved]}; it does not affect
the product structure.
\end{theorem}

The method is uniform and worth stating in one breath. Because
$a\circ b=a+\lambda_a(b)$, once the additive group $(G,+)$ and the homomorphism
$\lambda$ are fixed, the circle operation --- hence the factor set $\ovtau$ --- is
\emph{determined}, never solved for. The residue $[\nu]$ is the $D\to D$ block of
$\lambda$; the obstruction $\Om$ is the off-diagonal, section-varying part. The
block-triangular structure of the brace axioms then does all the work: the residue
is cut out by an action-only equation that never mentions the obstruction cochain,
and the obstruction is cut out by an affine equation whose homogeneous part is
residue-free. Orthogonality is not a coincidence of small examples; it is the
shape of the equations.

\medskip\noindent\textit{Outline.}
\S\ref{sec:omega-setup} records the factor-set calculus and the coupling equation.
\S\ref{sec:omega-identify} proves Theorem~\ref{thm:A-informal} (identification).
\S\ref{sec:omega-bilinear} proves Theorem~\ref{thm:B-informal} (bilinearity and
base-localisation). \S\ref{sec:omega-residue} isolates the residue $[\nu]$ as an
orthogonal factor and shows non-trivial kernels do not disturb $\Om$.
\S\ref{sec:omega-product} proves the product Theorem~\ref{thm:C-informal}.
\S\ref{sec:omega-corner} gives the genuine $|G|=16$ computational witness and
states the single open residual precisely. \S\ref{sec:omega-ledger} is the honest
ledger and the relation to Ferri's Remark~3.20.

%------------------------------------------------------------------------------
\section{Setup: factor sets and the coupling equation}\label{sec:omega-setup}

With $(G,+)$ fixed, a skew-brace structure making $D$ an ideal is \emph{exactly} a
homomorphism $\lambda\colon G\to\Aut(G,+)$ satisfying
$\lambda_{a\circ b}=\lambda_a\lambda_b$ (with $a\circ b:=a+\lambda_a(b)$) and
leaving $D$ invariant. The circle factor set is then read off from $\circ$; it is
never independently posited. This is the reduction that governs everything below:
the question ``does an extension with a prescribed residue exist?'' is a question
about $\lambda$, and the coordinates $\beta,\ovtau$ are the shadow of the single
identity $\lambda_{a\circ b}=\lambda_a\lambda_b$.

\begin{definition}[Inner group and residue]\label{def:residue}
The \emph{inner group} is $L:=\{\lambda_e\!\mid_D : e\in D\}=\operatorname{im}
(\lambda\!\mid_D)\le\Aut(D,+)$, with $L\cong(D,\circ)/\ker(\lambda\!\mid_D)$; write
$D_0:=\ker(\lambda\!\mid_D)$, so $D$ is a trivial brace iff $L=\{\id\}$ iff
$D_0=D$. The \emph{residue} of the section datum $\nu_h=\lambda_{s(h)}\!\mid_D$ is
the coset $[\nu_h]:=\nu_h L\in\Aut(D,+)/L$.
\end{definition}

In coordinates the brace axioms are equivalent to the following system (the
generalised Rathee--Yadav conditions; the additive-difference circle factor set is
written $T(h,k):=s(h)\circ s(k)-s(h\circ_H k)\in D$):
\begin{align}
  &\text{\rm(a)}\quad \mu_{h_1}\beta(h_2,h_3)+\beta(h_1,h_2{+}h_3)
      =\beta(h_1,h_2)+\beta(h_1{+}h_2,h_3), \tag{a}\\[2pt]
  &\text{\rm(b)}\quad T\ \text{is a circle $2$-cocycle (the $\circ$-associativity condition)}, \tag{b}\\[2pt]
  &\text{\rm(c)}\quad
     \mu_h\,\nu_h\,\beta(k,l)+T(h,k{+}_H l)
     = T(h,k)-\mu_{h\circ k}\mu_h^{-1}\beta(h,-h)+\beta(h\circ k,-h) \notag\\
  &\hphantom{\text{\rm(c)}\quad}\quad
     {}+\mu_{(h\circ k)-h}\,T(h,l)+\beta\bigl((h\circ k)-h,\ h\circ l\bigr), \tag{c}\\[2pt]
  &\text{\rm(coupling-1)}\quad \nu_h\,\mu_k\,\nu_h^{-1}=\mu_{\lambda^H_h(k)}. \tag{$\ast$}
\end{align}
Here $\lambda^H\colon(H,\circ_H)\to\Aut(H,+)$ is the induced $H$-brace map. Two
structural facts about this system carry the whole chapter, and we flag them now:
\begin{itemize}
\item In equations \textup{(a)--(c)}, the residue datum $\nu$ appears \emph{exactly
   once}, in the term $\mu_h\nu_h\beta(k,l)$ of \textup{(c)}.
\item Equation \textup{(coupling-1)} constrains the residue using only
   $\nu,\mu,\lambda^H$ --- it mentions neither $T$ nor $\beta$.
\end{itemize}
This is the block-triangularity: the residue side and the obstruction side of the
system decouple.

%------------------------------------------------------------------------------
\section{The obstruction is a skew-brace cohomology class}\label{sec:omega-identify}

We first pin $\Om$ down as a cohomology class and record what its vanishing means.
All of this section is \textsc{[peer-reviewed]} (referee: Rick); the cohomology
machinery is Rathee--Yadav's \cite{RatheeYadavExt}, invoked, not
re-derived.

Under (SH) the pair $(\beta,\ovtau)$, subject to (a), (b), (c), represents a class
\[
   \Om := [(\beta,\ovtau)] \in \Hcoh_{\Sb}(H,D)
     := Z^2_{\Sb}(H,D)/B^2_{\Sb}(H,D),
\]
where $Z^2_{\Sb}$ is the set of pairs satisfying the coupling (c) together with
$\beta\in Z^2_{\Gp}((H,+);\mu)$ and $\ovtau\in Z^2_{\Gp}((H,\circ);\sigma)$, and
$B^2_{\Sb}$ is generated by the section changes below. This is exactly the group
that classifies skew-brace extensions inducing the fixed triplet
$(\nu,\mu,\sigma)$. The circle datum $\ovtau$ and the additive-difference factor
set $T$ of \S\ref{sec:omega-setup} are one datum in two coordinate systems:
$\ovtau$ is the circle-multiplicative form and $T=\nu_{h\circ k}(\ovtau)$ its
additive image. We use $T$ in the cross-equation (c) --- where it must interact
additively with $\beta$ --- and $\ovtau$ for the circle cohomology class.

\begin{lemma}[Additive coboundary; $\beta=d\theta$]\label{lem:beta-cobound}
Under a section change $s'(h)=s(h)+\theta(h)$ with $\theta(0)=0$, the additive
factor set shifts by an additive coboundary and nothing else:
\[
   (\beta'-\beta)(h_1,h_2)=\mu_{h_2}\bigl(\theta(h_1)\bigr)-\theta(h_1{+}h_2)
   +\theta(h_2)=(d\theta)(h_1,h_2),
\]
independent of $\ovtau$, $\nu$ and $\sigma$.
\end{lemma}

This is the $\beta=d\theta$ convention of \cite{RatheeYadavExt}; we adopt it
verbatim. The lemma says the additive datum $\beta$ is exactly the group
$2$-cocycle of $0\to(D,+)\to(G,+)\to(H,+)\to0$ with $\mu$ its conjugation action;
so the forgetful assignment $\varphi_+\colon\Om\mapsto[\beta]\in
\Hcoh_{\Gp}((H,+),D;\mu)$ is well defined --- the coupling (c) constrains cocycles
but not coboundaries.

\begin{theorem}[Well-definedness of $\varphi_+$, \textsc{[peer-reviewed]}]
\label{thm:phiplus}
$\varphi_+$ is a well-defined homomorphism; the coupling (c) does not obstruct it.
The circle forgetful map $\varphi_\circ\colon\Om\mapsto[\ovtau]\in
\Hcoh_{\Gp}((H,\circ),D;\sigma)$ is defined symmetrically.
\end{theorem}

\begin{theorem}[Identification, \textsc{[peer-reviewed]}]\label{thm:identify}
Under (SH), for the extension $0\to D\to G\to H\to0$ with good triplet
$(\nu,\mu,\sigma)$,
\[
  \Om=0
  \iff \text{the sequence splits by a skew-brace section}
  \iff G\cong (G/D)\bic D
  \iff D \text{ has a sub-skew-brace complement.}
\]
No coprimality hypothesis is required (coprimality enters only Rathee--Yadav's
stronger uniqueness theorem). The first equivalence is intrinsic to the
Rathee--Yadav classification; the remaining ones are the registry result
\emph{corrected-equivalence-1-2-3}, via Ferri's Theorem~3.18 \cite{ferri-qsb}.
\end{theorem}

A \emph{sub-skew-brace complement} is a subset $K\subseteq G$ closed under both $+$
and $\circ$, with $K\cap D=\{0\}$ and $|K|\,|D|=|G|$; the equivalence records that a
one-sided (additive-only or circle-only) complement is strictly weaker, which is
precisely the content of the next section. The Lean corpus supplies machine-checked
\texttt{decide} certificates for the descent of $\varphi_+$ and for the
non-injectivity witness of \S\ref{sec:omega-bilinear} \textsc{[lean-verified]}.

%------------------------------------------------------------------------------
\section{Bilinearity: strictly finer than both Schreier classes}\label{sec:omega-bilinear}

The heart of the chapter's ``measurement'' claim is that $\Om$ is not the pair of
group classes $\bigl(\varphi_+(\Om),\varphi_\circ(\Om)\bigr)$; it is genuinely
bilinear, seeing a coupling between the two operations that neither operation sees
alone. Everything here is \textsc{[proved]}.

\begin{theorem}[The product map is not injective, \textsc{[proved]}]\label{thm:W4}
Let $\Phi:=(\varphi_\circ,\varphi_+)\colon
\Hcoh_{\Sb}(H,D)\to
\Hcoh_{\Gp}((H,\circ),D;\sigma)\times\Hcoh_{\Gp}((H,+),D;\mu)$.
Then $\ker\Phi\neq0$: there is a skew brace with $\Phi(\Om)=(0,0)$ yet $\Om\neq0$.
\end{theorem}

\begin{proof}[The $W_4$ witness]
Take additive group $(G,+)=\Z/2\times\Z/4$, multiplicative group
$(G,\circ)\cong\Z/4\times\Z/2$, and ideal $D=\{0\}\times\Z/4$ (a trivial-brace
ideal), with $H=G/D\cong\Z/2$. Then:
\begin{itemize}
\item the additive extension splits --- the additive complement
   $\{(0,0),(1,0)\}$ gives $[\beta]=\varphi_+(\Om)=0$;
\item the multiplicative extension splits --- the multiplicative complement
   $\{(0,0),(1,1)\}$ gives $[\ovtau]=\varphi_\circ(\Om)=0$;
\item but these two complements are \emph{distinct subsets}, and no single subset
   is a complement for both operations: there is no sub-skew-brace complement, so
   by Theorem~\ref{thm:identify}, $\Om\neq0$.
\end{itemize}
Hence $\Phi(\Om)=(0,0)$ with $\Om\neq0$. The obstruction lives in the coupling
term of (c); it is \emph{irreducibly bilinear}.
\end{proof}

The witness is exact and small enough to enumerate; a sweep over
$\Z/4,\ (\Z/2)^2,\ \Z/6,\ \Z/8,\ \Z/9,\ \Z/2\times\Z/4$ produces ten $W_4$-type
witnesses on $\Z/2\times\Z/4$, and locates a ``phantom'' family in which
$96$ of $144$ candidate ideals admit an ambient additive complement while $48$ do
not --- the additive picture alone cannot predict decomposition. Formally, the two
forgetful functors send $\Om$ to the additive and multiplicative Schreier classes
$[\omega_+]\in\Hcoh((G/D)_+;D_+)$ and $[\omega_\circ]\in\Hcoh((G/D)_\circ;D_\circ)$,
and $\Om$ may be nonzero while \emph{either or both} vanish. This is why the
decomposition axis of this chapter is strictly finer than the existence axis of
Chapter~\ref{ch:zs}: the group class $[\omega]$ is one leg of $\Phi$, and $\Phi$
has a kernel.

\paragraph{Base-localisation.}
In the groupoid (quiver skew brace) generalisation of the problem --- where $G$ is
a groupoid with a per-star additive structure and $D$ a wide normal sub-bundle ---
the obstruction is not merely finer, it is \emph{localised}. The naive
``heap-transport'' conjecture (that a decomposition of the vertex/loop skew brace
$L_\zeta$ suffices for a decomposition of $G$) is \emph{false}:

\begin{theorem}[Vertex splitting is necessary but not sufficient, \textsc{[proved]}]
\label{thm:heap-refute}
There is a connected quiver skew brace on $V\times\mathrm{Coarse}(\{0,1\})$ with
$V$ the Klein four-group and star group $(\Z/2)^3$, and a proper ideal sub-bundle
$D$, such that the vertex skew brace $L_\zeta$ decomposes ($D_\zeta$ has a
sub-skew-brace complement) yet $D$ has \emph{no} wide sub-quiver-skew-brace
complement, so $G\not\cong(G/D)\bic D$.
\end{theorem}

The correct invariant is a \emph{star complement}: $G\cong(G/D)\bic D$ iff there is
a multiplicative complement $K$ of $D_\zeta$ in $L_\zeta$ together with connecting
labels $(\tau_\mu)$ (with $\tau_\zeta=e$) such that
$A=\{(\zeta,\mu,k\tau_\mu):k\in K,\ \mu\in X\}$ is an additive complement of
$D_\zeta$ \emph{closed under $+_\zeta$ in the whole base star}. Closure on one base
star forces closure on every star, via the additive isomorphisms
$\rho_a=a\ract(-)$ carrying $D_\zeta$ to $D_\mu$; so the obstruction is base-$\zeta$
localised. For the Klein witness all single-operation shadows split --- additive
extension of the star, additive complement of $D_\zeta$, multiplicative complement
$K=\{0,3\}$, and the reduced base extension --- yet no star complement exists. This
is the same both-split-no-common-complement phenomenon as $W_4$, now with the extra
lesson that testing the loop shadow instead of the base star is precisely the
mistake the heap-transport conjecture made. (A brute-force sweep of the Klein
family found $504$ vertex-splitting-but-not-decomposing configurations and $0$ wide
complements; the star-transport lemma held in $1008/1008$ cases.)

%------------------------------------------------------------------------------
\section{The internal residue is an orthogonal factor}\label{sec:omega-residue}

We now separate $\Om$ from the datum most easily mistaken for part of it: the
residue $[\nu]$, which records how the internal action $\lambda$ relabels $D$ along
a section. All of this section is \textsc{[proved]}.

\begin{lemma}[Torsor lemma]\label{lem:torsor}
As $s'$ ranges over all normalised sections, $(\nu'_h)_{h\neq0}$ ranges over
exactly $T_\nu:=\prod_{h\neq0}\nu_h L$, a torsor under the free transitive action
of $L^{H\setminus0}$. Explicitly $\nu'_h=\nu_h\circ\lambda_{e(h)}\!\mid_D$ with
$e(h)=\lambda_{s(h)}^{-1}(\theta(h))$.
\end{lemma}

\begin{lemma}[Residue invariance]\label{lem:res-inv}
The cosets $[\nu_h]=\nu_h L\in\Aut(D,+)/L$ are section-independent, and
$\nu_{h_1}\nu_{h_2}=\nu_{h_1\circ h_2}\cdot\lambda_{\ovtau(h_1,h_2)}\!\mid_D$ with
$\lambda_{\ovtau}\!\mid_D\in L$. When $L\trianglelefteq\Aut(D,+)$ (automatic if
$\Aut(D,+)$ is abelian, e.g.\ $D$ cyclic), $[\nu]\colon(H,\circ)\to\Aut(D,+)/L$ is a
homomorphism.
\end{lemma}

\begin{lemma}[Obstruction $\nu$-freeness / descent, ``fact B'']\label{lem:nu-free}
With $(D,+),(D,\circ)$ abelian, $[\beta]$, $[\ovtau]$ and the predicate
``$\Om=0$'' are constant on the torsor $T_\nu$ and refer to no $\nu$. Concretely
the section-shifts $(\beta'-\beta)=d\theta$ and
$(\ovtau'-\ovtau)=d_\sigma e$ involve no $\nu$.
\end{lemma}

\begin{theorem}[Orthogonal direct factor, \textsc{[proved]}]\label{thm:ortho}
The section-independent invariants of the extension are
$\bigl(\mu,\sigma,[\nu],\Om\bigr)$, and the residue $[\nu]=([\nu_h])_h\in
\prod_h\Aut(D,+)/L$ is orthogonal to $\Om$:
\begin{enumerate}
\item its section ambiguity is a full torsor $T_\nu$ over $L^{H\setminus0}$ with
   invariant residue $[\nu]$ --- an ``$H^1$ over $L$'';
\item $[\beta]$, $[\ovtau]$ and splitting are $\nu$-free and constant on $T_\nu$,
   and $[\nu]$ is defined without reference to $\beta$, $\ovtau$ or the ambient
   defect.
\end{enumerate}
\end{theorem}

It is worth being exact about \emph{which} orthogonality holds, because a weaker
statement is false and the distinction is the crux.

\begin{proposition}[Not a gauge-orbit product --- the sharp negative, \textsc{[proved]}]
\label{prop:not-orbit}
The gauge action of the section group on the cochain triple
$(\nu,\beta,\ovtau)$ is \emph{not} a direct product of a $\nu$-action and an
obstruction-action: $\ker(\Psi_{\mathrm{obs}})+\ker(\Psi_\nu)\neq\Theta$. For the
braces $\mathtt{sb\#11},\mathtt{sb\#15}$ on $\Z/2\times\Z/4$ with $D=\{0\}\times\Z/4$
($L=\Aut(D,+)=\Z/2$), any section that flips $\nu_1$ forces
$d\theta(1,1)=(0,2)\neq0$.
\end{proposition}

That is: at the level of \emph{orbits} (raw cochains), moving the residue is
coupled to a nonzero change in the obstruction cochain. But that change is a
coboundary --- invisible to cohomology. Orthogonality holds for \emph{classes}, and
fails for orbits. The independence theorem is a statement about classes; the
negative proposition is the fence that stops us over-claiming it.

\paragraph{Non-trivial kernels do not disturb $\Om$.}
Dropping (SH) --- allowing $D$ to be a non-trivial brace --- changes the residue's
behaviour but not the obstruction:

\begin{theorem}[Kernel non-triviality is orthogonal, \textsc{[proved]}]\label{thm:kernel}
For an ideal $D$ with $(D,+)$ abelian, $\nu$ is section-independent for all
sections iff $\lambda_e\!\mid_D=\id$ for all $e\in D$ iff $D$ is a trivial brace.
When $D$ is a non-trivial brace, the sole effect is to make $\nu$
section-dependent; the descent maps $\varphi_+,\varphi_\circ$ and the ambient
common-complement defect are unchanged. Explicitly, for a lift $k$ of an order-$2$
quotient, $k\circ k=(k+k)+(\lambda_k-\id)(k)$, and $G\cong H\bic D$ iff there is
$k\in A:=\{k:k+k=0\}$ with $(\lambda_k-\id)(k)=0$; the defect
$(\lambda_k-\id)(k)$ uses $\lambda_k$ for a lift $k\notin D$, an ambient feature
into which $\nu$ never enters.
\end{theorem}

The witness makes the orthogonality vivid: on $\Z/2\times\Z/4$ with
$(D,+)\cong\Z/4$, the non-trivial-brace kernels $\mathtt{sb\#11},\mathtt{sb\#15}$
(with $(D,\circ)\cong\Z_2^2$) and the trivial-brace $\mathtt{sb\#1}$ share the same
$A=\{(1,0),(1,2)\}$, the same defect $(0,2)\neq0$, hence the same
$[\beta]=[\ovtau]=0$ and $\Om\neq0$; they differ \emph{only} in the number of
realisable $\nu_1$ ($1$ versus $2$). Kernel non-triviality is carried entirely by
$\nu$, and $\nu$ never touches $\Om$.

%------------------------------------------------------------------------------
\section{The product theorem}\label{sec:omega-product}

We can now upgrade orthogonality of a single pair of invariants to a global product
decomposition of the moduli of extensions. Fix the additive group $(G,+)$, the
subgroup $D$ (hence $\mu$ and the additive class $[\beta]$), the $D$-brace
$(L,\lambda\!\mid_D)$, the $H$-brace $(H,+,\circ_H)$, and the circle action
$\sigma$. Among all skew-brace structures $\circ$ on $(G,+)$ realising this data
with $D$ an ideal, taken up to equivalence of extensions, let $R$ be the set of
\emph{realised} residues $[\nu]$.

\begin{theorem}[Residue-independent obstruction group; product, \textsc{[proved]}]
\label{thm:product}
\leavevmode
\begin{enumerate}
\item For every $r\in R$, the equivalence classes of extensions with residue $r$
   form a torsor under an abelian group $\Omega^0$ that is \emph{independent of
   $r$}.
\item Consequently $\bigl([\nu],\Om\bigr)\colon\Ext\to R\times\Omega^0$ is a
   bijection: every realised residue co-occurs with every value of $\Om$.
\end{enumerate}
\end{theorem}

\begin{proof}[Proof (four steps)]
\emph{Step 1 (the residue side is $T$-free).} By (coupling-1), the admissible
residues are cut out using only $\nu,\mu,\lambda^H$ --- never $T$ or $\beta$. This
is block-triangularity: the $D\to D$ block of $\lambda$, carrying the residue,
satisfies a closed system not involving the off-diagonal obstruction block.

\emph{Step 2 (for fixed $\nu$, the valid $T$ form an $\varnothing$-or-torsor under
a $\nu$-free group $Z^0$).} Read (c) as an equation for $T$ with
$\beta,\mu,\nu,\lambda^H$ fixed. It is \emph{affine} in $T$, with homogeneous part
\[
   U(h,k{+}_H l)=U(h,k)+\mu_{(h\circ k)-h}\,U(h,l). \tag{$c^0$}
\]
The crux is that $(h\circ k)-h=\lambda^H_h(k)$ in $H$, so the twist
$\mu_{(h\circ k)-h}=\mu_{\lambda^H_h(k)}$ depends only on $\mu$ and $\lambda^H$
--- \emph{not on $\nu$}. Together with the homogeneous circle-cocycle condition
$(b^0)$ (also $\nu$-free), set $Z^0:=\{U:H\times H\to D \mid (b^0)\wedge(c^0)\}$,
manifestly independent of $\nu$ and $\beta$. If $T,T'$ both solve (b)$\wedge$(c)
then $U:=T-T'\in Z^0$ (the $\mu_h\nu_h\beta$ term and all $\beta$-terms cancel on
subtraction), and if $T$ solves (b)$\wedge$(c) and $U\in Z^0$ then $T+U$ does. So
the solution set $\mathcal T(\nu)$ is empty or a torsor under the $\nu$-free group
$Z^0$.

\emph{Step 3 (coboundaries are $\nu$-free; the obstruction group).} A change of
section shifts $T$ by a coboundary from a subgroup $B^0\subseteq Z^0$, shifts $\nu$
within its $L$-coset (residue preserved, Lemma~\ref{lem:res-inv}) and shifts
$\beta$ by an additive coboundary. The $T$-coboundaries $B^0$ are $\nu$-free
(Lemma~\ref{lem:nu-free}). Put $\Omega^0:=Z^0/B^0$. For each realised residue $r$,
the equivalence classes of extensions with residue $r$ are a torsor under
$\Omega^0$, and since $Z^0,B^0$ are $\nu$-free, $\Omega^0$ is residue-independent.
This is (1).

\emph{Step 4 (the product).} By (1), for each realised $r$ the class $\Om$ ranges
over the entire $\Omega^0$ as the extension ranges over the full torsor
$\Ext(r)$. Fixing a base-point identification $\Ext(r)\cong\Omega^0$, every value
of $\Omega^0$ is attained for every realised $r$, and because each $\Ext(r)$ is a
\emph{full} torsor the conclusion is base-point independent. Thus
$([\nu],\Om)$ is a bijection onto $R\times\Omega^0$.
\end{proof}

Nothing in the proof uses $[\beta]=0$: the theorem holds for every additive class.
The reader should note where each hypothesis is spent. That $(D,+),(D,\circ)$ are
abelian gives the factor-set calculus (a)--(c); fixing the $H$-brace $\lambda^H$
makes the twist in $(c^0)$ $\nu$-free; fixing $\mu,\sigma$ makes $(b^0)$ and $B^0$
$\nu$-free. This is the ``internal-replacement'' statement in its cleanest form:
\emph{changing the internal relabelling action $[\nu]$ does not change which
decomposition obstructions $\Om$ are available.} They are independent coordinates
of the moduli.

\begin{corollary}[Split $\perp$ residue, \textsc{[proved]}]\label{cor:split}
The split value $\Om=0$ is realisable iff $[\beta]=0$ (additive split), uniformly
in the residue: a sub-skew-brace complement is a fortiori an additive complement,
so $\Om=0\Rightarrow[\beta]=0$; conversely $[\beta]=0$ makes every admissible
residue realisable by a bicrossed product, which is split. Hence when $[\beta]=0$,
$\Ext\cong R\times\Omega^0$ with split and non-split each a residue-independent
subset; when $[\beta]\neq0$ no extension is split, and the realisable $\Om$-values
are the non-zero classes, attained uniformly across all realised residues.
\end{corollary}

%------------------------------------------------------------------------------
\section{A genuine $|G|=16$ witness, and the one open step}\label{sec:omega-corner}

The independence is easy to believe --- and easy to dismiss as a small-example
artifact. When $|D|=4$ one has $\Aut(\Z/4,+)=\Z/2$, so $L\neq1$ forces
$L=\Aut(D,+)$; the residue then has no room to be a \emph{proper} coset, and the
coupling never bites. A convincing witness must have $L\neq1$ \emph{and}
$L\subsetneq\Aut(D,+)$, which needs $|D|\ge8$, hence $|G|\ge16$. This corner is
where the phenomenon is real, and it is \textsc{[computed]}.

\begin{example}[The $|G|=16$ corner, \textsc{[computed]}]\label{ex:corner}
Take $(G,+)=\Z/2\times\Z/8$, $D=\{(0,x)\}\cong\Z/8$, $H=\Z/2$, so
$\Aut(D,+)=(\Z/8)^\times\cong\Z/2\times\Z/2$. Use the unit-scaling family
$\varphi_{u,t}(j,y)=(j,\ u\cdot y+j\cdot t)$ and the homomorphism
$U(i,x)=5^{x\bmod 2}\cdot 3^{i}\pmod 8$. Two structures differ only in a part
invisible on $D$ (so they share $L$ and $\nu$):
\[
   G_0:\ \lambda^0_{(i,x)}=\varphi_{U(i,x),\,0}
   \qquad
   G_1:\ \lambda^1_{(i,x)}=\varphi_{U(i,x),\,4i},
\]
where $T(i,x)=4i$ is a homomorphism $(G,\circ)\to\{0,4\}$. Then
$L=\{1,5\}$ is a \emph{proper, non-trivial} subgroup, $1\subsetneq L\subsetneq
\Aut(D,+)$; the base lift $s(1)=(1,0)$ gives $\nu_1=\ (\cdot 3)\notin L$ and residue
$[\nu_1]=\{3,7\}\neq L$. Both structures carry the \emph{same}
$(L,[\nu])=(\{1,5\},\{3,7\})$, yet $\Om=0$ for $G_0$ (the complement
$\{(0,0),(1,0)\}$ is closed under both operations) and $\Om\neq0$ for $G_1$
(every order-$2$ candidate fails $\circ$-closure: $(1,x)\circ(1,x)=(0,4)$). So
independence persists at genuine moving $\nu$ with non-trivial residue --- it is
not a $|D|=4$ size artifact.
\end{example}

The moving-$\nu$ cocycle equations of Example~\ref{ex:corner} are exactly
(a),(b),(c) of \S\ref{sec:omega-setup}, machine-checked on the witness; the
mechanism verdict is the one already used in the product theorem. At the
\emph{cochain} level there is a genuine coupling --- moving $\nu_h$ inside its
$L$-coset forces a compensating change in $T$. At the \emph{invariant} level there
is independence, because $a\circ b=a+\lambda_a(b)$ means that once $(G,+)$ (hence
$\mu,\beta$) and $\lambda$ (hence $[\nu]$ and the $L$-part) are fixed, $\circ$ and
therefore $T$ are \emph{determined}; equation (c) merely \emph{defines} $T$ from
$\beta$ up to a circle-coboundary.

\begin{table}[h]
\centering
\begin{tabular}{@{}llll@{}}
\toprule
regime & additive group & $[\beta]$ & result \\
\midrule
product enum & $\Z/2\times\Z/8$ & $0$ &
   every $\sigma$: full product of $([\nu],\Om)$; counts balanced \\
equiv-classes & $\Z/2\times\Z/8$ & $0$ &
   each admissible residue: exactly $2|L|$ classes ($|L|$ split, $|L|$ non-split) \\
equiv-classes & $\Z/16$ & $\neq0$ (ord $2$) &
   each residue: $2|L|$ classes, all non-split (uniform) \\
equiv-classes & $\Z/32$ & $\neq0$ (ord $4$) &
   residues $1,3,5,7$ all realised, each with exactly $4$ classes \\
engine-verify & $\Z/32$ & $\neq0$ &
   cross-equation (c) holds for all $28$ ideals ($H=\Z/4$) \\
\bottomrule
\end{tabular}
\caption{Exact enumeration corroborating Theorem~\ref{thm:product}. The constant
equivalence-class count per residue within each $(\sigma,L)$ block is the
prediction $|\Ext(r)|=|\Omega^0|$. The $[\beta]\neq0$, order-$4$ case ($\Z/32$)
realises all four admissible residues, refuting any fear that
$(\chi-1)_*[\beta]\neq0$ could block existence.}
\label{tab:enum}
\end{table}

\paragraph{The one open step, stated precisely.}
Theorem~\ref{thm:product} (the product on \emph{realised} residues) is
unconditional. What is \textsc{[computed]} but not \textsc{[proved]} is a single
surjectivity statement:

\begin{quote}
\emph{Surjectivity of the residue map for $[\beta]\neq0$:} every
coupling-admissible residue is actually realised (i.e.\ $R=R_{\mathrm{adm}}$) when
the additive extension is non-split.
\end{quote}

By Steps~2--3 this is equivalent to $\mathcal T(\nu)\neq\varnothing$ for every
admissible $\nu$, i.e.\ to the assertion that the RHS of (c) is always a
$T$-coboundary (equivalently: associativity of $\circ$ forces (c) to be solvable
for $T$, uniformly in the admissible residue). It is verified on $\Z/16$ and
$\Z/32$ (all admissible residues occur) and argued structurally on the witness,
but not closed in general. Three things should be emphasised, in the interest of
not over-claiming:
\begin{itemize}
\item the open residual concerns the \emph{residue map alone} --- it is not about
   the coupling term, and not about the product structure, both of which are
   established;
\item it does not affect Theorem~\ref{thm:product}, which is a statement about the
   realised set $R$;
\item the exact $(\sigma,L)$-block grouping is clean only for $H=\Z/2$ (where
   $\sigma$ is fully sampled at the single non-zero $h$); for $H=\Z/4$ the
   equivalence-class table is corroborated under a coarser $\sigma$-grouping, while
   the (c)-formula check there is exact.
\end{itemize}
Closing this residual is the business of a proof session, not this chapter; it is
flagged in the honest ledger below.

%------------------------------------------------------------------------------
\section{Honest ledger, and Ferri's Remark 3.20}\label{sec:omega-ledger}

\begin{center}
\begin{longtable}{@{}p{0.62\textwidth}l@{}}
\toprule
Statement & Grade \\
\midrule
\endhead
$\Om=[(\beta,\ovtau)]\in\Hcoh_{\Sb}(H,D)$ is the Rathee--Yadav class; $\varphi_+$
descends; $\Om=0\iff$ bicrossed (trivial-kernel SH) & \textsc{[peer-reviewed]} \\
Product map $\Phi=(\varphi_\circ,\varphi_+)$ non-injective ($W_4$) & \textsc{[peer-reviewed]} \\
$\Om$ is a bilinear class, strictly finer than both group Schreier classes;
base-$\zeta$ localised; heap-transport refuted & \textsc{[proved]} \\
$[\nu]$ is an orthogonal direct factor; gauge-orbit product is \emph{false},
orthogonality holds at the invariant level & \textsc{[proved]} \\
Non-trivial-brace kernels do not disturb $\Om$ (carried by $\nu$) & \textsc{[proved]} \\
Product theorem $\Ext\cong R\times\Omega^0$, $\Omega^0$ residue-independent
(on realised residues; unconditional) & \textsc{[proved]} \\
$|G|=16$ genuine-$L$ corner; moving-$\nu$ (c) machine-checked & \textsc{[computed]} \\
Surjectivity of the residue map for $[\beta]\neq0$ (every admissible residue
realised) & \textsc{[computed]} \\
$\varphi_+$ descent and $W_4$ non-injectivity \texttt{decide} certificates & \textsc{[lean-verified]} \\
\bottomrule
\end{longtable}
\end{center}

Two honesty notes are load-bearing and belong on the record.
\begin{itemize}
\item \emph{No socle-necessity claim.} An earlier draft of the identification
   asserted that separation ``requires $D\not\subseteq\Soc(G)$.'' On referee review
   (Rick) this was found to be false --- a trivial-kernel brace with $D=\Soc(G)$
   exactly nonetheless separates --- and the claim is \emph{deliberately omitted}
   here. It is not needed for any result in this chapter.
\item \emph{The general kernel is native.} The identification of $\Om$ with a
   published cohomology (Rathee--Yadav) is proved only for the trivial-kernel
   standing hypothesis. Housing the general non-trivial-kernel obstruction in a
   skew-brace $\Hcoh$ where $\nu$ is \emph{varying} cocycle data is native to this
   programme and not yet reduced to any published theory; we do not claim otherwise.
\end{itemize}

\paragraph{Relation to Ferri's Remark 3.20.}
Ferri \cite{ferri-qsb} asks, in the language of quiver skew braces, for a
group-theoretic reformulation of when an ideal is \emph{prunable} --- when the
quiver skew brace decomposes off that ideal. Theorem~\ref{thm:identify} supplies
exactly this: prunability of $D$ is the vanishing of the single skew-brace class
$\Om$, equivalently the existence of a sub-skew-brace (star) complement; and
Theorem~\ref{thm:B-informal} explains \emph{why} no purely group-theoretic
(single-operation) criterion can suffice --- the obstruction is bilinear, and lives
in the coupling that no forgetful functor to a group $\Hcoh$ can retain. The
base-localisation of Theorem~\ref{thm:heap-refute} is the groupoid refinement:
prunability is decided on one star and transported.

\paragraph{What this feeds.}
For the grant, the chapter closes the skew-brace / Yang--Baxter instance of the
decomposition programme with a clean separation of concerns: an \emph{existence}
axis ($[\omega]$, Chapter~\ref{ch:zs}) and a strictly finer \emph{decomposition}
axis ($\Om$, this chapter), the latter a bilinear invariant that is provably
independent of the internal relabelling residue. The independence is the
mathematical content of ``internal replacement'': one may re-index the internal
action of a composite system without changing which structural decompositions are
available to it. The single open surjectivity step is isolated, small, and does not
touch the product structure --- it is the natural target of the next proof session.
